% =====================================================================
%  Macros de los circuitos (CircuiTikZ ya se carga en el main con [american]).
%  Se puede \input en cualquier parte del documento antes de los circuitos.
% =====================================================================
\ctikzset{voltage dir=RP}

% ---- Etiquetas editables (cambiar con \renewcommand dentro de cada figura) ----
\providecommand{\etqFuente}{v_s}            % senal de entrada
\providecommand{\etqSec}{v_p}               % cada mitad del secundario (tab central)
\providecommand{\etqCarga}{R_L}             % resistencia de carga (reostato)
\providecommand{\etqSalida}{V_{out}}        % voltaje de salida
\providecommand{\etqCuno}{C_1}
\providecommand{\etqCdos}{C_2}
\providecommand{\etqRdummy}{1\,\Omega}      % valor de la R dummy
\providecommand{\sepCarga}{3.2}             % distancia R_L -> selector de C (subir si las etiquetas son largas)

% ---- Bloque: switch en paralelo con R dummy ----------------------------
% \dummyH{punto inicial}{largo}{switch}{resistencia}  horizontal, hacia la derecha
% \dummyV{punto inicial}{largo}{switch}{resistencia}  vertical, hacia abajo
\providecommand{\dummyH}[4]{%
  \draw (#1) to[short,*-] ++(0,0.6) to[nos, l={$#3$}] ++(#2,0) -- ++(0,-1.2)
             to[R, l={$#4$}] ++(-#2,0) -- ++(0,0.6);
  \draw (#1) ++(#2,0) node[circ]{};
}
\providecommand{\dummyV}[4]{%
  \draw (#1) to[short,*-] ++(-0.6,0) to[nos, l_={$#3$}] ++(0,-#2) -- ++(1.2,0)
             to[R, l_={$#4$}] ++(0,#2) -- ++(-0.6,0);
  \draw (#1) ++(0,-#2) node[circ]{};
}

% ---- Bloque: carga R_L con selector de capacitor en paralelo -----------
% \carga{nodo superior}{altura}{switch}
%   brazo al centro = abierto, izquierda = C1, derecha = C2
%   deja definida la coordenada (B) = nodo inferior de R_L
%   (flechas escritas como stealth-stealth para no chocar con babel-spanish)
\providecommand{\carga}[3]{%
  \coordinate (T) at (#1);
  \draw (T) to[vR, l_={$\etqCarga$}, v^={$\etqSalida$}, *-*] ++(0,-#2) coordinate(B);
  \draw (T) -- ++(\sepCarga,0) node[circ]{} coordinate(P) node[above]{$#3$};
  \draw[line width=0.8pt] (P) -- ++(0,-0.75);                          % brazo (abierto)
  \draw[stealth-stealth, thin] (P) ++(-125:0.55) arc[start angle=-125, end angle=-55, radius=0.55];
  \draw (P) ++(-0.7,-0.9) coordinate(K1) to[C, l_={$\etqCuno$}, o-*] (K1 |- B);
  \draw (P) ++( 0.7,-0.9) coordinate(K2) to[C, l={$\etqCdos$}, o-] (K2 |- B);
  \draw (B) -- (K2 |- B);
}

% ---- Transformador con tab central -------------------------------------
\providecommand{\trafoTab}{%
  \draw (0,0) to[sV, l={$\etqFuente$}] (0,4) -- (1.5,4) to[cute inductor] (1.5,0) -- (0,0);
  \draw[line width=1pt] (1.95,0.6) -- (1.95,3.4) (2.15,0.6) -- (2.15,3.4);
  \draw (2.6,4) to[cute inductor, mirror, l={$\etqSec$}] (2.6,2)
                to[cute inductor, mirror, l={$\etqSec$}] (2.6,0);
}
